\documentclass[11pt,a4paper]{article}
\usepackage[utf8]{inputenc}
\usepackage[margin=1in]{geometry}
\usepackage{array}
\usepackage{xcolor}
\usepackage{colortbl}
\usepackage{graphicx}

\definecolor{lightblue}{RGB}{230,240,250}

\title{Ejada Week 1 Task: OCI Infrastructure with Terraform}
\author{Mohamed Morsi}
\date{\today}

\begin{document}

\maketitle
\thispagestyle{empty}
\newpage

\section{Project Structure}

The project uses a flat file layout. Each file is named after the resource it configures:

\begin{itemize}
    \item \textbf{resources.tf} --- every resource block: VCN, Internet Gateway, Route Table, Security List, Subnet, Compute Instance, Reserved Public IP, Block Volume (+ attachment), File System (+ Mount Target and Export).
    \item \textbf{vcn.tf, internet\_gateway.tf, route\_table.tf, security\_list.tf, subnet.tf, compute\_instance.tf, block\_volume.tf, file\_system.tf} --- the \texttt{variable} declarations for each corresponding resource, kept in a file of the same name.
    \item \textbf{vnic\_data\_source.tf} --- a data source that looks up the private IP of the compute instance's VNIC, used to attach the reserved public IP.
    \item \textbf{output.tf} --- exposes the instance's public IP and the mount target's private IP after apply.
    \item \textbf{provider.tf} and \textbf{config.tf} --- see below.
    \item \textbf{terraform.tfvars} --- the actual values for every variable (credentials, CIDR blocks, display names, etc.). Not committed to version control.
\end{itemize}

\section{Main Configuration Files}

\subsection*{config.tf}
Declares the variables needed to authenticate against OCI: \texttt{user\_ocid}, \texttt{tenancy\_ocid}, \texttt{region}, \texttt{oci\_private\_key\_path}, and \texttt{fingerprint}.

\subsection*{provider.tf}
Configures the \texttt{oci} provider block, using API Key authentication with the variables declared in \texttt{config.tf}:

\begin{verbatim}
provider "oci" {
    auth = "APIKey"
    user_ocid = var.user_ocid
    tenancy_ocid = var.tenancy_ocid
    region = var.region
    private_key_path = var.oci_private_key_path
    fingerprint = var.fingerprint
}
\end{verbatim}

\section{Reserved Public IP}

By default, a compute instance's public IP is ephemeral, it changes if the instance is stopped/started or recreated. This project instead provisions a \textbf{reserved (static) public IP} via \texttt{oci\_core\_public\_ip}, with \texttt{lifetime = "RESERVED"}.

\begin{center}
\includegraphics[width=0.8\textwidth]{image.png}
\end{center}

\section{SSH Connection to Instance}

Connecting to the compute instance over SSH using its public IP address and the key pair configured in \texttt{instance\_ssh\_public\_key}.

\begin{center}
\includegraphics[width=0.8\textwidth]{1.png}
\end{center}

\section{Instance Reserved Public IP}

The reserved public IP address assigned to the instance, as shown in the OCI Console or via \texttt{terraform output week1\_public\_ip}.

\begin{center}
\includegraphics[width=0.8\textwidth]{image.png}
\end{center}

\section{Security List}

The security list (\texttt{Week 1 Security List}) allows inbound web traffic to the instance:

\begin{table}[h!]
\centering
\begin{tabular}{|l|l|l|l|l|}
\hline
\rowcolor{lightblue}
\textbf{Rule} & \textbf{Protocol} & \textbf{Source} & \textbf{Stateless} & \textbf{Destination Port} \\
\hline
Allow HTTP  & TCP (6) & 0.0.0.0/0 & No & 80 \\
\hline
Allow HTTPS & TCP (6) & 0.0.0.0/0 & No & 443 \\
\hline
\end{tabular}
\end{table}

\section{Route Table}

The route table (\texttt{Week 1 Route Table}) sends all outbound traffic to the Internet Gateway:

\begin{table}[h!]
\centering
\begin{tabular}{|l|l|l|}
\hline
\rowcolor{lightblue}
\textbf{Destination} & \textbf{Destination Type} & \textbf{Target} \\
\hline
0.0.0.0/0 & CIDR Block & Internet Gateway (Week 1 Public Gateway) \\
\hline
\end{tabular}
\end{table}

\section{Architecture Diagram}

The VCN (\texttt{Week 1 VCN}) is the network boundary for every resource in this project. An Internet Gateway is attached to the VCN and wired into the Route Table, which sends all outbound traffic (\texttt{0.0.0.0/0}) to that gateway --- this is what gives the subnet's resources a path to and from the internet.

Inside the VCN sits a single Public Subnet. That subnet uses the Route Table above for its outbound path, and is protected by the Security List, which only permits inbound TCP traffic on ports 80 and 443 (HTTP/HTTPS) from any source --- everything else is blocked.

Two resources run inside that subnet:
\begin{itemize}
    \item The \textbf{Compute Instance}, which has a Block Volume attached for extra disk space and a Reserved (static) Public IP so its address never changes across restarts.
    \item The \textbf{Mount Target}, which exposes a File System via an NFS Export, so the instance (or any other host inside the subnet) can mount shared storage over the network.
\end{itemize}

In short: traffic enters through the Internet Gateway, is routed into the subnet by the Route Table, is filtered down to ports 80/443 by the Security List, and finally reaches the Compute Instance --- which can read/write its own Block Volume and mount the shared File System through the Mount Target.

\begin{center}
\includegraphics[width=0.9\textwidth]{3.png}
\end{center}

\end{document}
